% ===========================================================================
\section{Dimension and Degree}\label{sec:dimdeg}
\warmth{0}\quad\whisper{Two numbers attached to an ideal, both squeezed out of one counting series.}

\begin{strand}
The two most important invariants of an ideal are its \keyword{dimension} and \keyword{degree}.
Both are read off a single generating function that \emph{counts monomials not in the ideal}, the
Hilbert series. The definition here is purely combinatorial; §\ref{sec:affine} reveals that it
matches the geometric dimension and degree of the variety.
\end{strand}

\block{Counting what is left}

\begin{definition}[Hilbert function \& series, Def.\ 1.20--1.21]\label{def:hilbert}
For a monomial ideal $I\subset K[\xb]$, the \keyword{Hilbert function} $h_I(q)$ is the number of
monomials of degree $q$ that are \fillin[3cm] $I$. Its generating function is the
\keyword{Hilbert series} $\HS_I(z)=\sum_{q\ge0} h_I(q)\,z^q$.
\end{definition}

\begin{yourturn}
Warm up by counting monomials directly.
\begin{itemize}[leftmargin=1.3em,itemsep=1pt]
  \item Zero ideal $I=\{0\}$: every monomial survives, so $h_{\{0\}}(q)=\binom{n+q-1}{n-1}$ and
        $\HS_{\{0\}}(z)=\fillin[2cm]$. \recall{This $h$ is a polynomial in $q$ of degree $n-1$.}
  \item Principal $I=\gen{\xb^{\mathbf a}}$ of degree $e$: subtract the multiples, giving
        $\HS_I(z)=\dfrac{1-z^{e}}{(1-z)^{n}}$. Read off: $\dim=\fillin[1.2cm]$, $\deg=\fillin[1.2cm]$.
\end{itemize}
\end{yourturn}

\begin{theorem}[shape of the series, Thm.\ 1.25]\label{thm:hs}
For any monomial ideal $I\subset K[\xb]$,
\[
  \HS_I(z)\;=\;\frac{\kappa_I(z)}{(1-z)^{n}},\qquad \kappa_I\in\Z[z],\ \kappa_I(0)=1,
\]
and there is a \keyword{Hilbert polynomial} $\HP_I$ with $\HP_I(q)=h_I(q)$ for all large $q$.
\end{theorem}
\begin{proof}
Inclusion--exclusion over the minimal generators $m_1,\dots,m_r$. For $\tau\subseteq\{1,\dots,r\}$
let $e_\tau=\deg\,\mathrm{lcm}(m_i:i\in\tau)$. Then $\kappa_I(z)=\sum_\tau(-1)^{|\tau|}z^{e_\tau}$.
\TODO{regroup to get $h_I(q)=\sum_\tau(-1)^{|\tau|}\binom{n+q-e_\tau-1}{n-1}$; check it is a polynomial for $q\gg0$}
\end{proof}

\block{Reading off dimension and degree}

\begin{definition}[dimension \& degree, Def.\ 1.28]\label{def:dimdeg}
Write $\HP_I(q)=\dfrac{g}{(d-1)!}\,q^{\,d-1}+(\text{lower order})$. Then the \keyword{dimension} of
$I$ is $d$ and the \keyword{degree} is the positive integer $g$. If $\HP_I\equiv0$ then $\dim=0$ and
$\deg I=\dim_K K[\xb]/I$ (a finite number).
\end{definition}

\recall{Slogan: dimension is the \emph{growth rate} of $h_I$; degree is its \emph{leading coefficient} (times $(d-1)!$).}

\begin{yourturn}
Two computations the book recommends.
\begin{itemize}[leftmargin=1.3em,itemsep=1pt]
  \item (Ex.\ 1.31) Principal ideal generated by a single monomial of degree $e$:
        confirm $\dim=n-1$ and $\deg=e$ from $\HP_I(q)=\dfrac{e}{(n-2)!}q^{\,n-2}+\cdots$.
  \item (Ex.\ 1.32) $n=2m$, $I=\gen{x_1x_2,\,x_3x_4,\,\dots,\,x_{2m-1}x_{2m}}$:
        find $\dim I=\fillin[1.2cm]$ and $\deg I=\fillin[1.2cm]$. \recall{Try $m=2$: $I=\gen{x_1x_2,x_3x_4}$. The degree is a power of two.}
\end{itemize}
\workspace[2]
\end{yourturn}

\begin{strand}
Why these two numbers (geometry preview, Remark 1.36 / Example 1.37). For the variety $\V(I)$:
$\dim=0$ \emph{iff} it is finitely many points, and then $\deg=$ the \emph{number} of points (counted
by the radical). A curve has $\dim 1$, a surface $\dim 2$. The degree measures how curvy the shape is.
A clean fact: a prime ideal has degree $1$ \emph{iff} it is generated by linear polynomials. For the
running cubic, $\gen f$ has $\dim 2,\deg 3$; its singular locus has $\dim 0,\deg 5$; together with $f$,
$\dim 0,\deg 4$ (the four nodes).
\end{strand}

\loose{Exercise 14 (book): let $X$ be a generic $2\times2$ matrix of variables, $I_s=\gen{\text{entries of }X^s}$. Investigate $\dim,\deg$ of $I_s$ for $s=2,3,4,\dots$. A small experiment with a CAS pays off.}
